% Part: counterfactuals
% Chapter: introduction
% Section: paradoxes-material

\documentclass[../../../include/open-logic-section]{subfiles}

\begin{document}

\olfileid{cnt}{int}{par}

\olsection{مادی شرطیے دے منطقی معمّے}

انگریزی دے ’جے \dots تاں \dots‘ والے بیاناں نوں $!A \lif !B$ نال لکھن دے
طریقے اُتے سبھ توں پہلاں اعتراض کرن والیاں وچ سی۔ آئی۔ لوئس سی۔
لوئس دا اعتراض وائٹ ہیڈ تے رسل دی \emph{Principia Mathematica} وچ مادی
شرطیے دی ورتوں اُتے سی۔ اوہ $\lif$ نوں ’لازم کردا اے‘ پڑھدے سن۔
لوئس نے ٹھیک اعتراض کیتا کہ جے $\lif$ دا مطلب ’لازم کردا اے‘ ہووے،
تاں ہر جھوٹا قضیہ~$p$ ایہ لازم کرے گا کہ $p$ توں~$q$ لازم آندا اے،
کیوں جے $p \lif (p \lif q)$ اُدوں سچ اے جدوں~$p$ جھوٹ اے۔ تے ہر سچا
قضیہ~$q$ ایہ لازم کرے گا کہ $p$ توں~$q$ لازم آندا اے، کیوں جے
$q \lif (p \lif q)$ اُدوں سچ اے جدوں $q$ سچ اے۔

منطق دے ماہر جان دے نیں کہ \emph{استلزام}، یعنی منطقی نتیجے دا تعلق،
!!{formula}s یا بیاناں وچکار اک تعلق اے، منطقی رابطہ نہیں۔ سو رَلا نہ
پئے، ایس لئی $\lif$ نوں ’لازم کردا اے‘ نہ پڑھیے۔\footnote{ریاضی دان تے
  کمپیوٹر سائنس دان ہجے وی ’$\lif$‘ نوں ’لازم کردا اے‘ عام پڑھدے نیں۔
  فلسفی لوئس دے دسے رَلے توں بچن لئی اینوں ’صرف جے‘ پڑھدے نیں۔}
جِنّی دیر اسیں ایویں نہیں پڑھدے، لوئس دی خاص اوکڑ پیدا ہی نہیں ہوندی۔
$p$، $q$ نوں ’لازم‘ نہیں کردا، بھانویں $p$ اک جھوٹے انگریزی جملے دا نشان
ہووے۔ $p \Entails q$ اے یا نہیں، ایہ طے کرن لئی \emph{ساریاں}
!!{valuation}s ویکھنیاں پیندیاں نیں۔ تے $p \Entails/ q$ اوہدوں وی اے
جدوں اسیں $p$ نوں ایسے جملے دا نشان بنائیے جیہڑا اتفاق نال جھوٹا اے۔

پر لوئس دی ہِیٹھلی مثال ورگے ’جے \dots تاں\dots‘ بیاناں وچ ہجے وی
اک عجیب گل اے:
\begin{quote}
جے چن کچّے پنیر دا بنیا اے، تاں $2+2=4$۔
\end{quote}
ایسی گل ہِیٹھلے استنباطاں وچ وی اے:
\begin{quote}
  چن کچّے پنیر دا نہیں بنیا۔ سو جے چن کچّے پنیر دا بنیا اے، تاں $2+2=4$۔

  $2+2 = 4$۔ سو جے چن کچّے پنیر دا بنیا اے، تاں $2+2=4$۔
\end{quote}
پر جے ’جے \dots تاں \dots‘ دا مطلب صرف $\lif$ ہووے، تاں ایہ جملہ
بغیر کسے اوکڑ دے سچ، تے ایہ استنباط بغیر کسے اوکڑ دے معتبر ہون گے۔

اک ہور مثال ہمیشہ سچے فارمولے $(!A \lif !B) \lor (!B \lif
!A)$ نال جڑی اے۔ ایہدے حساب نال جے اخبار وچوں دو خبری جملے~$S$ تے $T$
اٹکل نال چن لئیے، تاں ’جے $S$، تاں $T$، یا جے $T$، تاں~$S$‘ سچ
ہونا چاہیدا اے۔

\end{document}
